\documentclass[11pt,a4paper]{article}
\usepackage{amsmath,amssymb,amsthm,physics,geometry}
\usepackage{mathtools}
\usepackage{hyperref}
\geometry{margin=1in}

\title{\textbf{Exact Solution of the Repulsive Inverse--Square Potential\\
and Comparison with a Hard Wall at the Origin}}
\author{}
\date{}

\begin{document}
\maketitle

\begin{abstract}
We derive the exact radial solution of the Schrödinger equation
with repulsive inverse--square potential $V(r)=g/r^2$, $g>0$.
The regular solution is expressed in terms of Bessel functions
of non-integer order.
We analyze the near--origin behavior, prove essential self--adjointness,
and compute the scattering phase shift.
A precise comparison with the Dirichlet hard wall at $r=0$
is provided.
\end{abstract}

\section{Radial Equation}

Consider the three–dimensional Hamiltonian
\begin{equation}
H=-\frac{\hbar^2}{2m}\nabla^2+\frac{g}{r^2},
\qquad g>0.
\end{equation}

After separation of variables
\(
\psi(r,\theta,\phi)=R_l(r)Y_{lm}(\theta,\phi),
\)
define the reduced radial function
\(
u(r)=rR_l(r).
\)

The radial Schrödinger equation becomes
\begin{equation}
-\frac{\hbar^2}{2m}u''(r)
+
\frac{\hbar^2\lambda}{2m r^2}u(r)
=
Eu(r),
\end{equation}
where
\begin{equation}
\lambda
=
l(l+1)+\frac{2mg}{\hbar^2}.
\end{equation}

Define
\begin{equation}
\nu=\sqrt{\lambda+\frac{1}{4}}.
\end{equation}

For $g>0$ one has
\(
\nu>l+\tfrac{1}{2}.
\)

The equation reduces to
\begin{equation}
u''(r)+\left(k^2-\frac{\lambda}{r^2}\right)u(r)=0,
\qquad
k^2=\frac{2mE}{\hbar^2}.
\end{equation}

\section{Exact Scattering Solution ($E>0$)}

The general solution is
\begin{equation}
u(r)
=
\sqrt{r}
\left[
A J_\nu(kr)
+
B Y_\nu(kr)
\right].
\end{equation}

\subsection*{Behavior near the origin}

Using small–argument asymptotics,
\begin{align}
J_\nu(z) &\sim \frac{1}{\Gamma(\nu+1)}\left(\frac{z}{2}\right)^\nu,\\
Y_\nu(z) &\sim -\frac{\Gamma(\nu)}{\pi}\left(\frac{2}{z}\right)^\nu,
\end{align}
we obtain
\begin{equation}
u(r)\sim
A\, r^{\nu+\frac{1}{2}}
+
B\, r^{-\nu+\frac{1}{2}}.
\end{equation}

Since $\nu>1/2$ for $g>0$ and $l\ge0$,
the $Y_\nu$ branch diverges too strongly to be square–integrable.

Therefore physical regularity selects uniquely
\begin{equation}
u_{\text{reg}}(r)
=
C\sqrt{r}J_\nu(kr).
\end{equation}

Thus the Hamiltonian is essentially self–adjoint:
no boundary condition at $r=0$ is needed.

\section{Asymptotics and Phase Shift}

For large $r$,
\begin{equation}
J_\nu(kr)
\sim
\sqrt{\frac{2}{\pi kr}}
\cos\left(kr-\frac{\pi\nu}{2}-\frac{\pi}{4}\right).
\end{equation}

Hence
\begin{equation}
u(r)
\sim
\sqrt{\frac{2}{\pi k}}
\cos\left(kr-\frac{\pi\nu}{2}-\frac{\pi}{4}\right).
\end{equation}

Compare with free radial solution (no $g$):
\begin{equation}
u^{(0)}(r)
=
\sqrt{r}J_{l+\frac{1}{2}}(kr)
\sim
\sqrt{\frac{2}{\pi k}}
\cos\left(
kr-\frac{\pi}{2}\left(l+\frac{1}{2}\right)-\frac{\pi}{4}
\right).
\end{equation}

Therefore the phase shift is
\begin{equation}
\boxed{
\delta_l
=
-\frac{\pi}{2}
\left(
\nu - \left(l+\frac{1}{2}\right)
\right).
}
\end{equation}

Since $g>0$ implies $\nu>l+\tfrac{1}{2}$,
the phase shift is negative: purely repulsive scattering.

\section{Comparison with Hard Wall at $r=0$}

Consider instead a free particle with Dirichlet boundary
\begin{equation}
u(0)=0.
\end{equation}

For the free Hamiltonian,
the regular solution is
\begin{equation}
u^{\text{wall}}(r)=\sqrt{r}J_{l+\frac{1}{2}}(kr).
\end{equation}

Near the origin,
\begin{equation}
u^{\text{wall}}(r)\sim r^{l+1}.
\end{equation}

For inverse–square repulsion,
\begin{equation}
u(r)\sim r^{\nu+\frac{1}{2}}.
\end{equation}

Thus the repulsive inverse–square potential
produces stronger suppression:
\begin{equation}
\nu+\frac{1}{2}
>
l+1.
\end{equation}

\subsection*{Interpretation}

\begin{itemize}
\item Hard wall enforces $u(0)=0$ but leaves power fixed by $l$.
\item Repulsive $g/r^2$ modifies the effective centrifugal barrier:
\begin{equation}
l(l+1)\longrightarrow l(l+1)+\frac{2mg}{\hbar^2}.
\end{equation}
\item The wavefunction vanishes with higher power.
\item No reflection at a sharp boundary occurs.
\item The effect is purely phase–shifted scattering.
\end{itemize}

\section{Effective Angular Momentum}

Define
\begin{equation}
l_{\text{eff}}
=
\nu-\frac{1}{2}.
\end{equation}

Then the solution equals exactly the free radial solution
with noninteger angular momentum:
\begin{equation}
u(r)
=
\sqrt{r}J_{l_{\text{eff}}+\frac{1}{2}}(kr).
\end{equation}

Hence:

\begin{center}
\textbf{Repulsive inverse--square potential
is equivalent to shifting $l\to l_{\text{eff}}$.}
\end{center}

This is stronger than a hard wall,
because it modifies the centrifugal barrier rather than imposing
a boundary constraint.

\section{Bound States}

For $g>0$, there are no bound states.

For $E<0$, the solution involves modified Bessel functions:
\begin{equation}
u(r)=\sqrt{r}K_\nu(\kappa r),
\end{equation}
but these are not normalizable at infinity without additional
confinement.

\section{Conclusion}

For $g>0$:

\begin{enumerate}
\item Exact solution: $u(r)=\sqrt{r}J_\nu(kr)$.
\item Hamiltonian is essentially self–adjoint.
\item No bound states.
\item Scattering phase shift:
\(
\delta_l=-\frac{\pi}{2}(\nu-l-\tfrac{1}{2}).
\)
\item Physically equivalent to fractional angular momentum.
\item Stronger suppression than a hard wall at the origin.
\end{enumerate}

Thus the repulsive inverse–square potential
acts not as a boundary,
but as a dynamical renormalization of angular momentum.
\end{document}
